\chapter{第 5 章分层答案}
\label{app:ch05-solutions}

\section*{口述概念}
\begin{enumerate}[leftmargin=*]
  \item 偏导是沿坐标基向量的方向导数；方向导数逐方向给一个标量或向量极限。Fréchet 可微要求存在同一个线性映射 $A$，使余项为 $o(\lVert h\rVert)$。标量函数用列梯度表示该线性映射，向量函数用 Jacobian 表示。正文反例沿坐标轴偏导为零，却沿 $y=x^2$ 跳到 $1/2$，说明偏导存在甚至不保证连续。
  \item 本文向量和梯度为列向量，$J_g$ 的行对应输出、列对应输入，所以 $dg=J_gdx$。外层标量函数满足 $dh=\nabla h^Tdg=\nabla h^TJ_gdx$；与 $dh=\nabla(h\circ g)^Tdx$ 比较后转置整个系数，得到 $\nabla(h\circ g)=J_g^T\nabla h$。矩阵变量梯度由 $df=\operatorname{tr}((\nabla_Xf)^TdX)$ 定义。
  \item 若混合二阶偏导在邻域连续，Clairaut 定理保证 Hessian 对称；只在一点存在不够。解析推导能给一般公式和假设；自动微分能精确求执行图的局部导数，但不能验证目标语义；有限差分是独立数值检查，却受步长、舍入和不可微点影响。三者互补，不能互相替代。
\end{enumerate}

\section*{笔试推导}
\begin{enumerate}[leftmargin=*]
  \item 写 $g(x+u)=g(x)+Au+r_g(u)$，其中 $\lVert r_g(u)\rVert=o(\lVert u\rVert)$；再写 $h(y+v)=h(y)+Bv+r_h(v)$。取 $v=Au+r_g(u)=O(\lVert u\rVert)$，则
  \[
   h(g(x+u))-h(g(x))=BAu+B r_g(u)+r_h(v).
  \]
  后两项除以 $\lVert u\rVert$ 趋零，故导数为 $BA$。若 $h$ 标量值，$J_g^T\nabla h$ 的形状是 $(d\times m)(m\times1)=d\times1$。
  \item
  \[
   d(x^TAx)=dx^TAx+x^TAdx=x^T(A^T+A)dx.
  \]
  因此 $\nabla f=((A+A^T)/2)x+b$，$H_f=(A+A^T)/2$。二阶 Taylor 主项为 $\tfrac12h^TH_fh$；反对称部分在该二次型中消失。
  \item 令 $r=X\beta-y$，则 $dr=X\,d\beta$，故 $d\ell=r^TX\,d\beta=[X^Tr]^Td\beta$，梯度为 $X^T(X\beta-y)$。由 $X^{-1}X=I$ 微分，$(dX^{-1})X+X^{-1}dX=0$，左右整理得 $dX^{-1}=-X^{-1}(dX)X^{-1}$。Jacobi 公式给 $d\log\det X=\operatorname{tr}(X^{-1}dX)$，与矩阵梯度定义比较得 $\nabla_X\log\det X=X^{-T}$。
\end{enumerate}

\section*{数值编程}
\begin{enumerate}[leftmargin=*]
  \item 必须取 $(A+A^T)x/2+b$，本例为 $(0,-0.5)^T$；二次型精确算术下中心差分无截断项，极小步长主要受舍入影响。
  \item 真导数为 1，中心差分误差主项为 $h^2/6$，很小时浮点消去可能使误差反弹。
  \item 左、右导数为 $-1,1$，导数不存在；中心差分恒为零，不能证明可微。
\end{enumerate}

\section*{研究判断}
\begin{enumerate}[leftmargin=*]
  \item 先固定最小标量目标、参数点和 dtype；列出每个中间量形状及 Jacobian 约定；读取独立解析闭式结果。再比较自动微分与解析参照，检查 detach、原地修改、广播、转置和不可微算子；最后扫描中心差分步长并比较左右差商。每次只改变一个因素，解释误差随步长的变化。
\end{enumerate}

\section*{基础衔接：独立计算}
$\nabla f=(x+y-1)(1,1)^T$，Hessian 为全 1 的 $2\times2$ 矩阵。极小点是直线 $x+y=1$，目标为零。沿 $(1,-1)$ 的曲率为零，Hessian 半正定但不正定。
